\section{Related Work}
\label{appendix:related_work}

\subsection{Transformers in Reinforcement Learning}
According to the survey by \citet{li2023survey}, Transformers \citep{vaswani2017attention} are increasingly utilized in reinforcement learning (RL) for various tasks, including representation learning of individual observations and their histories, as well as model learning, as seen in Dreamer by \citet{hafner2019dream}. In our research, we focus on the application of Transformers in sequential decision-making and in developing generalist agents. Same as AD \citep{laskin2022context}, Headless-AD also utilizes transformers as the base model for our approach.

The incorporation of transformers as a sequence modeling tool in RL began with the Decision Transformer (DT) by \citep{hu2022transforming}, which is trained autoregressively on offline datasets of state, action, and return-to-go tuples. Unlike conventional RL, which focuses on return maximization, DT generates appropriate actions during inference by conditioning on specified return-to-go values. The Trajectory Transformer by \citet{janner2021offline} is an alternative approach using beam search to bias trajectory samples based on future cumulative rewards. Building on this, the Multi-Game Decision Transformer (MGDT) by \citet{lee2022multi} improves upon DT with enhanced transfer learning capabilities for new games, eliminating the need for manual return specification. Similarly, the Switch Trajectory Transformer by \citep{lin2022switch} expands on the Trajectory Transformer to facilitate multi-task training. However, when transitioning to novel tasks, these approaches require fine-tuning the model.

% Further, Hindsight Information Matching (HIM) \citep{furuta2021generalized} suggests a unified framework that enables conditioning on any hindsight information instead of return-to-go. Choosing different hindsight information helps to solve stochastic environments as in ESPER \citep{paster2022you} and DoC \citep{yang2022dichotomy} or augmenting a dataset as in Q-learning Decision Transformer (QDT)\citep{yamagata2023q}.

\subsection{Offline Meta-RL}
Traditional RL agents, tailored to specific environments, struggle with novel tasks. In contrast, by training across a diverse array of tasks, Meta-RL equips agents with adaptable exploration strategies and an understanding of common environmental patterns \citep{duan2016rl, wang2016learning, rajeswaran2017towards, zhang2018dissection, team2021open, nikulin2023xlandminigrid}. Offline Meta-RL, a subset of this approach, trains agents solely on pre-existing datasets, without direct environmental interaction. A critical challenge here is MDP ambiguity, where task-specific policies misinterpret data due to dataset biases. \citet{dorfman2021offline} propose a data collection method to mitigate this issue, and suggest an approach that treats Meta-RL as a Bayesian RL problem for optimal exploration in new tasks. \citet{li2020focal} introduce FOCAL, a framework separating task identification from control. However, their method relies on a strict mapping assumption that fails in certain scenarios, such as those with sparse rewards. \citet{li2021provably} enhance FOCAL with an attention mechanism and improved metric learning, showing greater robustness in scenarios with sparse rewards and domain shifts. However, most Offline Meta-RL methods rely on explicit task modeling, which can introduce limiting biases. Alternatively, In-Context Learning implicitly infers tasks from environmental interactions, offering a potentially more flexible approach.


\subsection{In-Context Learning in RL}
In-Context Learning (ICL) is an ability of a pretrained model to adapt and perform a new task given a context with examples $\{x_i, f(x_i)\}^n$, where $f$ is a function that gives ground truth targets \cite{brown2020language, wang2023images}. Previous work on ICL \cite{mirchandani2023large} showed that Large Language Models operate as General Pattern Machines and are able to complete, transform and improve token sequences, even when the sequences consist of randomly sampled tokens. This ability supports our design choice to randomly encode the actions of agents. The research on Transformer Circuits \cite{elhage2021mathematical} gives evidence of Transformers' ability to copy tokens, either literally or on a more abstract level, explaining their ICL capabilities. This is particularly useful for our model, which explicitly incentivizes an abstract copying of the action embeddings.

Efforts to blend In-Context Learning (ICL) abilities of transformers with reinforcement learning (RL) are gaining traction, promising to create adaptable RL agents for real-world scenarios with varying conditions. Attempts to transfer ICL features to RL include Prompt-Based Decision Transformers \citep{xu2022prompting}, which leverage task-specific demonstration datasets as prompts, exhibiting strong few-shot learning without weight updates. \citet{laskin2022context} introduced Algorithm Distillation (AD) that trains a policy improvement operator using data from agent-environment interactions of learning RL algorithm, predicting the next action autoregressively. A key strategy here is across-episode training for capturing policy improvement. \citet{lee2023supervised} developed Decision Pretrained Transformer (DPT) using supervised training with unordered environmental interactions as context to predict optimal actions. Under certain conditions, DPT can be proved to implement posterior sampling, resulting in near-optimal exploration. However, DPT's reliance on an optimal policy during training, not always available, is a limitation, similar to AD's constraint on action space structure, restricting their use as fully generalist agents. \citet{lin2023transformers} provided theoretical analysis of AD and DPT, offering guarantees about learning from base algorithms. \citet{wang2024transformers} also provided a theoretical analysis showing that transformers may implement several RL algorithm in-context. \citet{zisman2023emergence} propose a method for distilling the improvement operator from a demonstrator whose actions are initially noisy, with the noise level decreasing progressively throughout the data collection phase. This approach simplifies the data generation compared to AD as it does not require logging the training process of RL models and relaxes DPT's limitation because the demonstrator may be suboptimal. \citet{kirsch2023towards} demonstrated that data augmentation, through random projection of observations and actions, enhances task distribution and generalization on new domains, boosting ICL capabilities. \citet{raparthy2023generalization} study the effect of model size and dataset properties on the success of in-context learning. Other research, acknowledging Transformers' limitations when it comes to long sequences, explores alternative models such as S4 \citep{gu2021efficiently}. \citet{lu2023structured} adapted S4 for RL, showing that it is capable of surpassing LSTM in performance and Transformers in runtime, indicating potential in RL ICL applications.

In contrast to previous works that tackle changing reward distributions or environment dynamics as changing goals, we expand the application of ICLRL to changing action spaces. Our use of random embeddings, motivated by the preparation of the model to unseen actions, may also improve ICL abilities by data augmentation, as suggested by \citet{kirsch2023towards}. 

\subsection{Discarding the Linear Layer}
The Headless LLM approach, introduced by \citet{godey2023headless}, utilizes Contrastive Learning to eliminate the language head from the model architecture. Rather than generating a probability distribution over tokens, it directly predicts token embeddings. This strategy aims to enhance training and inference speeds by discarding a substantial linear layer. While in the RL context, where action spaces are typically small, this might not significantly impact runtime, it does offer a key advantage, as the model is no longer dependent on the number of actions and can therefore handle variable-length action spaces. Unlike Headless LLM, our approach does not rely on contrastive loss for learning action representations. Instead, we use fixed action embeddings, and the model is tasked to output the embedding of the next action. During inference, an action is chosen from a categorical distribution, where the logits are the similarities between the predicted embedding and available actions. Thus, the role of contrastive loss is to enhance the likelihood of selecting the correct action while diminishing the probabilities of other actions. 

Wolpertinger by \citet{dulac2015deep} also employs a similar approach of removing the linear head to improve the train and inference speeds. The authors suggest associating each action with an embedding and performing the training in a continuous action space. A specific action index is chosen as a nearest neighbor of the predicted embedding, and its effect is treated as environment dynamics. 
Most importantly, this nearest neighbor selection process can be optimized using an approximate algorithm. This leads to logarithmic time complexity, which is particularly advantageous in environments with large action spaces. Thus, we believe that Headless-AD is also capable of performing well on large action spaces, benefiting from runtime gains of approximate nearest neighbor lookup. However, Wolpertinger's usage of fixed action embeddings prevents introduction of new actions, highlighting the significance of Headless-AD's usage of randomized action embeddings.